\chapter{Inheritance and Polymorphism}

\section{Introduction}
In session 6, we bundled our code into a single class: \texttt{IrisRuleClassifier}.
It had two main jobs:
\begin{enumerate}
    \item \textbf{Decide:} Does the sample have \texttt{petal\_length < threshold}? (This is specific to this rule).
    \item \textbf{Evaluate:} Loop through the dataset and calculate accuracy. (This is general).
\end{enumerate}

So far, having one class is already an improvement over passing a \texttt{settings} dictionary everywhere.
But a new design problem appears as soon as you want \textbf{more than one type} of classifier.

\noindent
What if we want to create a new classifier, such as a \texttt{NearestCentroidClassifier} (a different prediction strategy) or even a simple \texttt{RandomGuessClassifier}?
With the session 6 design, every new class would need:
\begin{itemize}\setlength{\itemsep}{2pt}
    \item its own \texttt{predict(...)} logic (different for each classifier type),
    \item and a copy of the same \texttt{evaluate(...)} loop (identical across classifier types).
\end{itemize}

\noindent
\textbf{Why is that a problem?}
\begin{itemize}\setlength{\itemsep}{2pt}
    \item Copy-paste creates duplication. Duplication increases bugs.
    \item If you fix a mistake in \texttt{evaluate()}, you must remember to fix it in every class.
    \item Adding a new classifier becomes ``rewrite everything again'' instead of ``write only what is different''.
\end{itemize}

This session (Session 7), we introduce \textbf{Inheritance}.
We will extract the shared logic (evaluation) into a \textbf{Parent Class} (Base), and keep only the specific logic (prediction rules) in the \textbf{Child Classes}.
\noindent
\textbf{Why introduce inheritance now?} Because you now have enough experience to see the repeating pattern:
many classifiers differ in \texttt{predict()}, but they can share \texttt{evaluate()}.
Inheritance is one clean way to reuse that shared behavior.

\subsection{Transition Map: session 6 Class $\rightarrow$ session 7 Inheritance}

\input{session7/tab7.tex}

\subsection{Deliverable}
Open or create \texttt{session7/session7\_iris\_template.py}.
By the end of this session, you will have:
\begin{itemize}\setlength{\itemsep}{2pt}
    \item a parent class \texttt{ClassifierBase} that defines the shared evaluation loop,
    \item two child classes that implement different \texttt{predict()} rules,
    \item a \texttt{main()} script that can evaluate multiple classifier objects in the same way.
\end{itemize}
\noindent
The key outcome is that you can extend the system by writing only what is new, without duplicating the evaluation code.

\section{The Parent Class }


\subsection{Step 1: The Parent Class (\texttt{ClassifierBase})}
This class will hold the logic that is \textit{shared} by all classifiers.
For us, that is the `evaluate` method. It doesn't know \textit{how} to predict yet, but it knows how to \textit{score} predictions.

\begin{minted}[fontsize=\small, linenos]{python}
class ClassifierBase:
    def __init__(self):
        pass

    def predict(self, sample):
        # The parent doesn't know how to predict!
        # Child classes MUST replace this method with their own logic.
        raise NotImplementedError("Child class must implement predict()")

    def evaluate(self, dataset):
        # This logic is exactly the same as session 6!
        # It calls self.predict(), which will be provided by the child class.
        correct = 0
        total = 0
        
        for sample in dataset:
            prediction = self.predict(sample)
            
            # Simple binary truth check
            truth = "setosa" if sample["species"] == "setosa" else "not_setosa"
            
            if prediction == truth:
                correct += 1
            total += 1
            
        return correct / total if total > 0 else 0.0
\end{minted}

\noindent
\textbf{Task:} Copy this parent class into your file.
\noindent
\textbf{Before you code (reasoning):} decide what should be shared across classifiers.
Evaluation is shared (loop, counting, accuracy). Prediction is not shared (rules differ), so the parent class should \textit{not} try to implement it.

\noindent
\textbf{Why does \texttt{predict} raise an error?}
This is intentional.
It is a beginner-friendly way to say: ``If you try to use \texttt{ClassifierBase} directly, Python should stop you, because the base class is incomplete.''
The error also helps you notice quickly if you forgot to implement \texttt{predict()} in a child class.

\noindent
\textbf{After you code (verification):} do not call \texttt{ClassifierBase().evaluate(...)} yet.
That should fail, and it should fail \textit{for a clear reason}: \texttt{predict()} is not implemented.


\section{Inheritance and Child Classes}


\subsection{Step 2: The First Child (\texttt{RuleClassifier})}
Now we recreate our session 6 classifier, but this time it \textbf{inherits} from \texttt{ClassifierBase}.
It only needs to provide the \texttt{\_\_init\_\_} and \texttt{predict} logic. It gets \texttt{evaluate} for free!

\begin{minted}[fontsize=\small, linenos]{python}
class RuleClassifier(ClassifierBase):  # <--- Inherits from ClassifierBase
    def __init__(self, threshold=2.0):
        super().__init__()  # Initialize the parent
        self.threshold = threshold

    def predict(self, sample):
        # Specific logic for THIS classifier
        if sample["petal_length"] < self.threshold:
            return "setosa"
        return "not_setosa"
\end{minted}

\noindent
\textbf{Task:}
\begin{enumerate}
    \item Define this class below \texttt{ClassifierBase}.
    \item Create an instance: \texttt{rule\_clf = RuleClassifier(2.0)}
    \item Call \texttt{print(rule\_clf.evaluate(dataset))}
\end{enumerate}
\noindent
\textbf{Before you code (reasoning):} your child class should contain only what is \textit{special} about the rule-based classifier.
The evaluation loop is not special, so it stays in the parent.

\noindent
\textbf{After you code (verification):} notice that you can call \texttt{rule\_clf.evaluate(dataset)} even though you did not write \texttt{evaluate} inside \texttt{RuleClassifier}.
That is the practical benefit of inheritance: reuse shared code without copy-paste.

\noindent
\textbf{Common beginner pitfall:} forgetting to inherit.
If you accidentally write \texttt{class RuleClassifier:} (missing \texttt{(ClassifierBase)}), you will not get \texttt{evaluate()} for free.

\subsection{Step 3: The Second Child (\texttt{NearestCentroidClassifier})}
Now we see the power of inheritance. We can add a completely different classifier without rewriting the evaluation loop.
This classifier classifies a sample by comparing its \texttt{petal\_length} to two target values and picking the closer one.
It is not ``better'' than the rule classifier; it is just \textbf{different}, and it demonstrates extensibility.

\begin{minted}[fontsize=\small, linenos]{python}
class NearestCentroidClassifier(ClassifierBase):
    def __init__(self, setosa_target=1.4, non_setosa_target=4.5):
        super().__init__()
        self.setosa_target = setosa_target     # Target petal_length for setosa
        self.non_setosa_target = non_setosa_target # Target for others

    def predict(self, sample):
        # Logic: Which target value is the sample's petal_length closer to?
        val = sample["petal_length"]
        dist_to_setosa = abs(val - self.setosa_target)
        dist_to_non = abs(val - self.non_setosa_target)

        if dist_to_setosa < dist_to_non:
            return "setosa"
        return "not_setosa"
\end{minted}

\noindent
\textbf{Task:}
\begin{enumerate}
    \item Define this class.
    \item Create an instance: \texttt{centroid\_clf = NearestCentroidClassifier()}
    \item Call \texttt{print(centroid\_clf.evaluate(dataset))}
\end{enumerate}
\noindent
\textbf{Before you code (reasoning):} identify what must be provided by the child class.
The child must provide \texttt{predict()}, and it must return the same label strings that \texttt{evaluate()} expects (here: \texttt{"setosa"} or \texttt{"not\_setosa"}).

\noindent
\textbf{After you code (verification):} you should be able to evaluate this classifier on the same dataset without writing any new evaluation code.


\section{Polymorphism}


\subsection{Step 4: Putting it all together}
We can now loop through a list of \textit{different types} of classifiers and treat them all the same because they share the same Parent Interface.

\begin{minted}[fontsize=\small, linenos]{python}
# 1. Define dataset
dataset = [
    {"petal_length": 1.4, "species": "setosa"},
    {"petal_length": 4.5, "species": "versicolor"},
    {"petal_length": 6.0, "species": "virginica"},
]

# 2. Create a list of different classifiers
classifiers = [
    RuleClassifier(threshold=2.0),
    NearestCentroidClassifier(setosa_target=1.5, non_setosa_target=5.0),
    RuleClassifier(threshold=1.5),
]

# 3. Evaluate them all in one loop
print("=== Model Comparison ===")
for clf in classifiers:
    acc = clf.evaluate(dataset)
    # Uses the class name to identify which one ran
    print(f"{clf.__class__.__name__}: Accuracy = {acc}")
\end{minted}
\noindent
\textbf{What is happening conceptually?}
\begin{itemize}\setlength{\itemsep}{2pt}
    \item Each object has a different \texttt{predict()} implementation.
    \item They all share the same \texttt{evaluate()} implementation from the parent class.
    \item The loop does not need to know the exact class type. It only needs to know that the object supports \texttt{evaluate()}.
\end{itemize}
\noindent
This ``treat different objects the same way'' behavior is a key OOP idea called \textbf{polymorphism}.

\section{Why is this better than session 6?}
In session 6, we only had one type of classifier.
In session 7, you build an extensible system where:
\begin{itemize}\setlength{\itemsep}{2pt}
    \item shared logic lives in one place (\texttt{ClassifierBase.evaluate}),
    \item new classifier types only implement what is unique (\texttt{predict} and any new attributes),
    \item your evaluation pipeline automatically works for any new child class, as long as it follows the same interface.
\end{itemize}
\noindent
This is a refactoring mindset you will reuse often: extract shared code into a reusable component, then specialize only the parts that differ.

\subsection{Submission Requirement}
Add your student label at the end:
\begin{minted}[fontsize=\small, linenos]{python}
print("=== Student Label ===")
print("ID: ...")
\end{minted}
Run the script passing the dataset to both classifiers.
